\clearpage\thispagestyle{empty}
\pdfbookmark[0]{Quellen zur Pallas-Anzeige von 1810}{anzeige1810-sources}
\begin{center}\textbf{Quellen zur Pallas-Anzeige von 1810}\end{center}
\noindent Diese KI-Übertragung des deutschen Originals in heutiges Deutsch folgt dem Text in Gauß’ \textit{Werken}, Band~VI (Göttingen, 1874), S.~61--64. Die Anzeige erschien in den \textit{Göttingischen gelehrten Anzeigen}, 198.~Stück, 13.~Dezember 1810, S.~1969--1973. \href{https://www.digitale-sammlungen.de/de/view/bsb10538678?page=939}{Erstdruck in der Bayerischen Staatsbibliothek}. \href{https://gdz.sub.uni-goettingen.de/id/PPN236007467}{Historischer deutscher Text in den Werken}.
\par\smallskip\noindent Die Anzeige bespricht die lateinische Abhandlung über die elliptischen Elemente der Pallas, die am 25.~November 1810 vorgelegt wurde. Sie enthält eigene Tabellen mit vier Systemen von Elementen und sechs Paaren von Abweichungen.
\par\smallskip\noindent Übertragung und Satz wurden mit KI erstellt.
\clearpage
